\item Considere la velocidad con que camina un animal bípedo o cuadrúpedo, modele una pierna que no está contactando a la tierra como una barra uniforme de longitud $\ell$, que se hace pivotar como un péndulo físico a través de una mitad de un ciclo, en resonancia. Sea que $\theta_{\text{máx}}$ represente su amplitud.
\begin{enumerate}
    \item Demuestre que la rapidez del animal está dada por la expresión:
    $$ v = \sqrt{\frac{6g\ell}{\pi}} \sin\theta_{\text{máx}} $$
    si $\theta_{\text{máx}}$ es lo suficientemente pequeño como para que el movimiento sea casi armónico simple. Una relación empírica que se basa en el mismo modelo y se aplica en un gran rango de ángulos es:
    $$ v = \sqrt{\frac{6g\ell}{\pi}} \cos(\theta_{\text{máx}}/2) \sin\theta_{\text{máx}} $$
    \item Evalúe la rapidez con que camina un ser humano con una pierna de longitud $0.850\text{ m}$ y una amplitud de balanceo de pierna de $28.0^\circ$. ¿Qué longitud de pierna daría el doble de la rapidez para la misma amplitud angular?
\end{enumerate}